\section{Theta--Mellin reciprocity and conservation of the polar structure}
\label{rec:espectral}

\subsection{Periodization of the generated lattice}

The lattice \(\Lambda_t\) of Theorem~\ref{rec:thm:red} has an analytic reading that preserves its dependence on the register. For \(x>0\), define
\begin{equation}
 \Theta_t(x)=\sum_{v\in\Lambda}
       \exp\!\left(-\pi_{\HMT}x\|G_tv\|^2\right).
 \label{rec:eq:theta}
\end{equation}
In this realization, the previously generated coordinate \(\pi_{\HMT}\) is identified with the circular normalization of Fourier analysis. That identification follows the construction of the lattice and the flow. To simplify the analytic notation, we shall use \(\pi\) with this meaning.

On every compact interval of \(t\), the singular values of \(G_t\) are bounded above and below by positive constants. The lattice has \(O(R^{24})\) points in a ball of radius \(R\). Comparison over annuli with a Gaussian proves absolute and uniform convergence of \eqref{rec:eq:theta} and its derivatives when \(x\) remains in a compact subset of \((0,\infty)\). The mode \(v=0\) contributes exactly one.

\begin{theorem}[Reciprocity of the theta reading]
\label{rec:thm:theta}
For every \(x>0\) and \(t\in\RR\),
\begin{equation}
 \Theta_t(x)=x^{-12}\Theta_{-t}(1/x).
 \label{rec:eq:theta-reciproca}
\end{equation}
\end{theorem}
\begin{proof}
With the convention \(\widehat f(\xi)=\int f(y)e^{-2\pi i\langle y,\xi\rangle}\,dy\), the Gaussian \(f_x(y)=e^{-\pi x\|y\|^2}\) in dimension twenty-four has transform
\(\widehat f_x(\xi)=x^{-12}e^{-\pi\|\xi\|^2/x}\).
Periodize over \(\Lambda_t\). Integrating the periodization over a fundamental domain shows that the Fourier coefficient indexed by \(\xi\in\Lambda_t^*\) is \(\widehat f_x(\xi)/\operatorname{covol}(\Lambda_t)\). Both Gaussian series converge absolutely; evaluating the expansion at the origin gives
\[
 \sum_{v\in\Lambda_t}f_x(v)
 =\operatorname{covol}(\Lambda_t)^{-1}
   \sum_{\xi\in\Lambda_t^*}\widehat f_x(\xi).
\]
Substitute \(\operatorname{covol}(\Lambda_t)=1\) and \(\Lambda_t^*=\Lambda_{-t}\). The proof retains the constant mode on both sides.
\end{proof}

The power \(x^{-12}\) comes from rank twenty-four. Inversion of the parameter follows from the dual lattice; it is not added as a symmetry independent of the reading.

\subsection{Meromorphic continuation and functional equation}

For \(\Re s>12\), let
\begin{equation}
 \mathcal E_t(s)=\sum_{v\in\Lambda\setminus\{0\}}\|G_tv\|^{-2s},
 \qquad \mathcal L_t(s)=\pi^{-s}\Gamma(s)\mathcal E_t(s).
 \label{rec:eq:epstein}
\end{equation}
The lattice-point count above proves absolute convergence in that half-plane. Applying the gamma integral term by term yields
\[
 \mathcal L_t(s)=\int_0^\infty(\Theta_t(x)-1)x^{s-1}\,dx.
\]

\begin{theorem}[Polar decomposition and bilateral reciprocity]
\label{rec:thm:mellin}
The function \(\mathcal L_t\) extends meromorphically through
\begin{align}
 \mathcal L_t(s)={}&\frac1{s-12}-\frac1s\nonumber\\
 &+\int_1^\infty\left[(\Theta_t(x)-1)x^{s-1}
                 +(\Theta_{-t}(x)-1)x^{11-s}\right]\,dx.
 \label{rec:eq:mellin-partida}
\end{align}
Its only poles are simple, at \(s=0,12\), with residues \(-1,+1\), respectively. It satisfies
\begin{equation}
 \mathcal L_t(s)=\mathcal L_{-t}(12-s),\qquad
 \Xi_t(s):=s(s-12)\mathcal L_t(s)=\Xi_{-t}(12-s),
 \label{rec:eq:funcional}
\end{equation}
and \(\Xi_t\) is entire.
\end{theorem}
\begin{proof}
Split the integral at \(x=1\). For \(0<x<1\), theta reciprocity gives
\[
 \Theta_t(x)-1=x^{-12}-1+x^{-12}(\Theta_{-t}(1/x)-1).
\]
The first two terms integrate to \(1/(s-12)-1/s\). In the third, the change of variable \(y=1/x\) gives
\(\int_1^\infty(\Theta_{-t}(y)-1)y^{11-s}\,dy\).
The positive minimum length of the lattice ensures exponential decay of \(\Theta_t(x)-1\) as \(x\to\infty\). On every compact set of \(s\), that decay dominates every power of \(x\) and \(\log x\); the integrals and their complex derivatives therefore converge uniformly and define entire functions. The residues are read from the rational part. Simultaneously replacing \((t,s)\) by \((-t,12-s)\) exchanges the integrals and leaves the rational part invariant. Multiplication by \(s(s-12)\) removes both poles.
\end{proof}

The centre \(s=6\) follows from the rank and the exponent \(-2s\) adopted in \eqref{rec:eq:epstein}. This collection is the Epstein completion of the deformed lattices. Its spectral variable and normalization remain distinct from those of the Riemann zeta function; the relation proved between the two constructions is the reciprocity of their circular and lattice realizations.

\begin{corollary}[Polar conservation and antisymmetric component]
\label{rec:cor:polar}
The functions
\[
 \mathcal L_t(s)-\mathcal L_0(s),\qquad
 \mathcal A_t(s)=\frac{\mathcal L_t(s)-\mathcal L_{-t}(s)}2
\]
are entire. Moreover,
\[
 \mathcal A_t(12-s)=-\mathcal A_t(s),\qquad \mathcal A_t(6)=0.
\]
\end{corollary}
\begin{proof}
The rational part of \eqref{rec:eq:mellin-partida} is independent of \(t\), so it cancels in both differences. The functional equation changes the sign of the second. Evaluation at its fixed point gives \(2\mathcal A_t(6)=0\).
\end{proof}

The deformation therefore resides in an entire difference and preserves the polar data exactly. The central zero belongs to the bilateral difference, without imposing a zero on each \(\mathcal L_t\). The poles described are singularities of a Mellin transform; the physical connection with the mass operator is established later through the action section, with its own variable and reader.

\subsection{Stable reconstruction from the bilateral archive}
\label{rec:archivo}

For the unitary cyclic shift \(S\), the preceding recovery law provides
\[
 T_j=\frac{8I+S^j}{9},\qquad D_j=\frac{\sqrt8(S^j-I)}9,
 \quad j=3,4,
\]
\begin{equation}
 m=(D_3K,D_4T_3K)=MK,\qquad LM=P_{11},\qquad\|L\|=\frac94.
 \label{rec:eq:archivo}
\end{equation}
Consequently,
\begin{equation}
 u=P_3Lm.
 \label{rec:eq:recuperacion-u}
\end{equation}
With the incidence chart, the marked lattice, and the regional ratio \(\rho\) fixed, this identity recovers successively \(u/\|u\|\), \(g_t\), \(G_t\), \(\Lambda_t\), \(\Theta_t\), and \(\mathcal E_t\). Reconstructing this collection requires neither the terminal \(T_4T_3K\) nor the register's uniform charge. The geometric antecedents remain among the fixed data.

\begin{theorem}[Geometric and spectral stability of recovery]
\label{rec:thm:estabilidad}
Suppose that \(m'=m+e\), \(\|e\|\le\varepsilon\), and \(9\varepsilon/4<\|u\|\). Let \(u'=P_3Lm'\), and reconstruct the flows with \(u'/\|u'\|\), leaving \(\rho,P_3,L\) and the charts unchanged. If \(a=|t\log\rho|\), then
\begin{equation}
 \|g'_t-g_t\|,\ \|G'_t-G_t\|
 \le\frac92\,a e^a\frac{\varepsilon}{\|u\|}.
 \label{rec:eq:estabilidad}
\end{equation}
With \(H_K=\operatorname{diag}(h_i)\), \(b=|t|\|H'_K-H_K\|\), we have
\begin{align}
 b&\le\frac{9a\varepsilon}{2\|u\|},\nonumber\\
 \Theta_t(e^{2b}x)&\le\Theta'_t(x)\le\Theta_t(e^{-2b}x),\label{rec:eq:theta-error}\\
 e^{-2\sigma b}\mathcal E_t(\sigma)&\le\mathcal E'_t(\sigma)
             \le e^{2\sigma b}\mathcal E_t(\sigma),\qquad\sigma>12.
 \label{rec:eq:epstein-error}
\end{align}
\end{theorem}
\begin{proof}
The norm of the recovery operator gives \(\|u'-u\|\le9\varepsilon/4\), ensuring \(u'\ne0\). The triangle inequality and the reverse triangle inequality for norms imply
\[
 \left\|\frac{u'}{\|u'\|}-\frac{u}{\|u\|}\right\|
 \le\frac{2\|u'-u\|}{\|u\|}.
\]
Every exponent lies between \(-a\) and \(a\). Applying the mean value theorem to the exponential and taking the diagonal operator norm proves \eqref{rec:eq:estabilidad}; repeating each entry twice does not change that norm. The same estimate before exponentiation gives the bound on \(b\).

Since both generators are diagonal in the same chart,
\[
 e^{-2b}\|G_tv\|^2\le\|G'_tv\|^2\le e^{2b}\|G_tv\|^2.
\]
Applying the decreasing exponential and summing proves \eqref{rec:eq:theta-error}. Raising to \(-\sigma\) and summing within the domain of convergence proves \eqref{rec:eq:epstein-error}. The bounds apply to all vectors, not to a finite truncation of the series.
\end{proof}

The error \(\varepsilon\) is measured in the normalized memories of \eqref{rec:eq:archivo}. If the rational records \(z=m/\sqrt8\) are retained, their error is multiplied by \(\sqrt8\) before applying the estimate. Independent errors in the regional ratio or the charts require their own bounds. The result recovers a collection of geometries and analytic readings, with quantitative control of the archive that determines it.
